\section{信托法}
\begin{enumerate}
  \item 1 道题。
  \item 基本同民法的代理。
  \item 信用委托的理财方式。需要转移所有权。
  \item 委托人基于对受托人的信任，将其财产权委托给受托人，由受托人按委托人的意愿以自己的名义，为受益人的利益或\zd{特定目的}，进行管理和处分的行为。
  \item 分类。
        \begin{enumerate}
          \item 民事信托、营业信托（信托公司）、公益信托。
          \item 自益 \yjt 受益人就自己一个、他益。
          \item 撤销、不可撤销。
          \item 多个委托人。单一、集合。
        \end{enumerate}
  \item 设立。3 确定。信托财产确定、受益人确定、信托目的确定。
  \item 成立。合同；遗嘱。
  \item 无效。不能设立讨债和诉讼目的的信托。
  \item 撤销。
  \item 信托财产范围。原始 + 收益 + 限制流通财产（经审批，如贵金属）。
  \item 独立性。独立委托人、受托人、受益人、委托人（之后）的债权人。
  \item 原则上不得强制执行。
        \begin{enumerate}
          \item 设立之前，债权人已享有优先受偿权。
          \item 处理信托产生的正常债务。
          \item 信托财产本身负担的税费。
        \end{enumerate}
  \item 委托人。完民。
        \begin{enumerate}
          \item 对受托人处分行为的撤销权
                \begin{enumerate}
                  \item 找法院；恢复原状，予以赔偿；主观起算 1 年。
                  \item 违反目的处分财产，处理信托事务不当。
                \end{enumerate}
          \item 解任权。重大过失。
          \item 变更受益人、处分信托受益权。受益人侵权行为；同意。基本同民法丧失继承资格情况。
          \item 解除权。自益自己就可以解除；受益人重大侵权；同意。
        \end{enumerate}
  \item 受托人。完民。
        \begin{enumerate}
          \item 可以是受益人，但不能是唯一的受益人。通谋。
          \item 原则亲自处理信托事务，有约从约。同民法代理。
          \item 分别管理、分别记账；公平对待；禁止自我交易（一手托两家）。
          \item 职责终止不意味着信托终止 \yjt 换人 \yjt 有约从约 \yjt 委托人指定。
        \end{enumerate}
  \item 受益人。
        \begin{enumerate}
          \item 信托受益权；均等比例；可处分。
          \item 都放弃 \yjt 信托终止。
          \item 解任权、撤销权同受托人，冲突找法院，一个人及于全部。
        \end{enumerate}
  \item 终止。连续性。
\end{enumerate}
